\chapter{XÂY DỰNG HỆ THỐNG VÀ TÍCH HỢP TRỢ LÝ AI}

\section{Tổ chức thư mục mã nguồn}

Dự án được chia thành 3 thư mục chính gồm \texttt{backend/} (API quản lý và cổng MCP), \texttt{ai-runner/} (dịch vụ chạy lệnh gọi AI) và \texttt{frontend/} (giao diện React), được tóm tắt trong Bảng~\ref{tab:source-tree}.

\begin{table}[!htbp]
\centering
\footnotesize
\renewcommand{\arraystretch}{1.05}
\caption{Cấu trúc các thư mục chính trong kho mã nguồn}
\label{tab:source-tree}
\begin{tabularx}{\textwidth}{|L{4.2cm}|Y|L{3.1cm}|}
\hline
\thead{Thư mục} & \thead{Chức năng chính} & \thead{Công nghệ}\\
\hline
\texttt{backend/app/core} & Đọc biến môi trường (\texttt{config.py}), kết nối CSDL (\texttt{database.py}), mã lỗi (\texttt{errors.py}) và tạo/kiểm tra JWT (\texttt{security.py}). & Python 3.12, Passlib, PyJWT\\
\hline
\texttt{backend/alembic} \& \texttt{backend/seed} & Quản lý phiên bản bảng CSDL (Alembic) và kịch bản nạp dữ liệu mẫu (\texttt{seed\_demo.py}). & SQLAlchemy 2.0, Alembic\\
\hline
\texttt{backend/app/models} \& \texttt{schemas} & Khai báo 14 bảng dữ liệu và các lớp kiểm tra dữ liệu đầu vào/đầu ra. & SQLAlchemy 2.0, Pydantic v2\\
\hline
\texttt{backend/app/api/v1} & 10 nhóm API: \texttt{auth}, \texttt{users}, \texttt{channels}, \texttt{campaigns}, \texttt{contents}, \texttt{schedules}, \texttt{metrics}, \texttt{stats}, \texttt{brand}, \texttt{notifications}. & FastAPI Routers\\
\hline
\texttt{backend/app/assistant} & Máy chủ công cụ FastMCP (\texttt{mcp\_server.py}, \texttt{tools.py}), luồng chat SSE (\texttt{chat\_service.py}) và xử lý xác nhận thẻ đề xuất (\texttt{actions\_service.py}). & FastMCP, HTTPX, SSE\\
\hline
\texttt{ai-runner/app} & Dịch vụ nhận yêu cầu gọi AI (\texttt{main.py}, \texttt{runner.py}), xếp hàng tuần tự và chạy tiến trình \texttt{opencode run}. & FastAPI, Asyncio, OpenCode CLI\\
\hline
\texttt{frontend/src/features} & 8 nhóm màn hình: \texttt{auth}, \texttt{dashboard}, \texttt{campaigns}, \texttt{contents}, \texttt{schedules}, \texttt{metrics}, \texttt{assistant}, \texttt{settings}. & React 18, TypeScript, TanStack Query\\
\hline
\texttt{frontend/src/components} & Khung giao diện (\texttt{AppShell.tsx}, \texttt{Sidebar.tsx}), ngăn kéo chat AI (\texttt{AssistantDrawer.tsx}), thẻ đề xuất (\texttt{ProposedActionCard.tsx}) và chọn model (\texttt{ModelSelector.tsx}). & Tailwind CSS v4, ECharts, FullCalendar\\
\hline
\end{tabularx}
\end{table}

\section{Xây dựng các chức năng quản lý}

\subsection{Phân quyền theo vai trò và người tạo chiến dịch}

Khi người dùng gửi yêu cầu lên API, hàm \texttt{get\_current\_user} trong \texttt{backend/app/api/deps.py} kiểm tra chữ ký JWT, trạng thái hoạt động (\texttt{is\_active == True}) và đối chiếu \texttt{token\_version} với CSDL. Bên cạnh việc phân quyền theo 2 vai trò \texttt{MANAGER} và \texttt{MARKETER}, hệ thống kiểm tra thêm cột \texttt{owner\_id} trên từng chiến dịch để đảm bảo nhân viên này không thể tự ý sửa hoặc xóa chiến dịch, bài viết của nhân viên khác (Bảng~\ref{tab:scope-matrix}).

\begin{table}[!htbp]
\centering
\footnotesize
\renewcommand{\arraystretch}{1.05}
\caption{Bảng phân quyền thao tác theo vai trò và người phụ trách chiến dịch}
\label{tab:scope-matrix}
\begin{tabularx}{\textwidth}{|L{2.2cm}|L{2.7cm}|L{3.1cm}|Y|C{2.0cm}|}
\hline
\thead{Dữ liệu} & \thead{Thao tác} & \thead{Quản lý (\texttt{MANAGER})} & \thead{Nhân viên phụ trách (\texttt{owner\_id})} & \thead{Nhân viên khác}\\
\hline
\multirow{2}{2.2cm}{\raggedright Chiến dịch} & Xem danh sách / Chi tiết & Được phép & Được xem để phối hợp & Được xem\\
\cline{2-5}
& Sửa / Xóa / Chia ngân sách & Được phép & Được sửa/xóa chiến dịch của mình & Chặn (HTTP 403)\\
\hline
\multirow{3}{2.2cm}{\raggedright Bài viết} & Soạn / Sửa / Gửi duyệt & Được phép & Thực hiện trên chiến dịch của mình & Chặn (HTTP 403)\\
\cline{2-5}
& Duyệt / Từ chối bài & Duyệt bài của người khác (\texttt{CANNOT\_}\allowbreak\texttt{REVIEW\_OWN}) & Không có quyền duyệt (HTTP 403) & Chặn (HTTP 403)\\
\cline{2-5}
& Xóa bài viết & Cấm xóa bài đã duyệt (\texttt{CONTENT\_LOCKED}) & Chỉ xóa khi ở \texttt{DRAFT}/\texttt{REJECTED} & Chặn (HTTP 403)\\
\hline
Lịch \& Chỉ số & Lên lịch / Nhập số liệu & Được phép & Thực hiện trên chiến dịch của mình & Chặn (HTTP 403)\\
\hline
Nhân sự \& Kênh & Quản lý Users / Brand Kit & Toàn quyền quản trị & Chỉ được xem thông tin thương hiệu \& kênh & Chặn (HTTP 403)\\
\hline
\end{tabularx}
\end{table}

\subsection{Nhập số liệu quảng cáo (tay \& CSV 2 bước) và tính chỉ số hiệu quả}

Người dùng có thể nhập số liệu quảng cáo theo ngày bằng biểu mẫu hoặc tải lên file báo cáo CSV qua 2 bước:
\begin{enumerate}[noitemsep,topsep=2pt]
 \item \textbf{Bước 1 (Xem trước tại \texttt{POST /api/v1/metrics/import/preview})}: Máy chủ đọc từng dòng trong file CSV, kiểm tra tên kênh có thuộc chiến dịch hay không, đồng thời kiểm tra điều kiện số liệu hợp lý $\text{views} \ge \text{clicks} \ge \text{conversions} \ge 0$ và $\text{cost} \ge 0$. Hệ thống trả về danh sách dòng hợp lệ và danh sách dòng lỗi kèm lý do cụ thể trên màn hình.
 \item \textbf{Bước 2 (Xác nhận lưu tại \texttt{POST /api/v1/metrics/import/commit})}: Chỉ những dòng hợp lệ được người dùng xác nhận mới được ghi vào bảng \texttt{metrics}.
\end{enumerate}

Để tránh việc tính nhẩm sai hoặc lỗi chia cho 0 khi chiến dịch mới tạo chưa có lượt xem/lượt nhấp, mô-đun \texttt{backend/app/api/v1/stats.py} và hàm MCP \texttt{get\_stats} tính sẵn các chỉ số bằng hàm \texttt{\_safe\_div(num, den)} trước khi trả về cho giao diện và Trợ lý AI:
\begin{align}
\text{CTR (\%)} &= \begin{cases} \dfrac{\sum \text{clicks}}{\sum \text{views}} \times 100 & \text{nếu } \sum \text{views} > 0 \\ 0,0 & \text{nếu } \sum \text{views} = 0 \end{cases} \\
\text{CVR (\%)} &= \begin{cases} \dfrac{\sum \text{conversions}}{\sum \text{clicks}} \times 100 & \text{nếu } \sum \text{clicks} > 0 \\ 0,0 & \text{nếu } \sum \text{clicks} = 0 \end{cases} \\
\text{CPC} &= \begin{cases} \dfrac{\sum \text{cost}}{\sum \text{clicks}} & \text{nếu } \sum \text{clicks} > 0 \\ 0,0 & \text{nếu } \sum \text{clicks} = 0 \end{cases} \\
\text{CPA} &= \begin{cases} \dfrac{\sum \text{cost}}{\sum \text{conversions}} & \text{nếu } \sum \text{conversions} > 0 \\ 0,0 & \text{nếu } \sum \text{conversions} = 0 \end{cases}
\end{align}
Chỉ số $\text{CPM} = (\sum \text{cost} / \sum \text{views}) \times 1000$ và phần trăm ngân sách đã chi tiêu $\text{used\_pct} = (\sum \text{cost} / \text{budget}) \times 100\%$ cũng được tính theo cùng quy tắc trên.

\subsection{Giao diện thực tế của hệ thống}

Giao diện người dùng được chia thành 3 khu vực như Hình~\ref{fig:ui-wireframe}: thanh menu bên trái (\texttt{Sidebar.tsx}) gồm các mục \textit{Tổng quan}, \textit{Chiến dịch}, \textit{Nội dung}, \textit{Duyệt nội dung}, \textit{Lịch đăng}, \textit{Trợ lý AI} và \textit{Cài đặt}. Khu vực nội dung chính ở giữa hiển thị bảng số liệu, biểu đồ, hàng đợi duyệt bài (có nhãn cảnh báo \textit{``Bạn không thể tự duyệt nội dung do chính mình tạo''}) và lịch đăng. Bên phải hoặc tại trang \texttt{/assistant} là khung Trợ lý AI kết nối trực tiếp với chiến dịch đang chọn.

\begin{figure}[!htbp]
\centering
\setlength{\tabcolsep}{4pt}
\begin{tabular}{cc}
 \fbox{\includegraphics[width=0.47\textwidth]{screenshots/dashboard-1280px.png}} &
 \fbox{\includegraphics[width=0.47\textwidth]{screenshots/campaign-detail-1280px.png}} \\[2pt]
 {\scriptsize (a) Trang Tổng quan Marketing và các thẻ KPI thực tế} &
 {\scriptsize (b) Trang Chi tiết chiến dịch, mục tiêu và phân bổ ngân sách kênh} \\[6pt]
 \fbox{\includegraphics[width=0.47\textwidth]{screenshots/reviews-1280px.png}} &
 \fbox{\includegraphics[width=0.47\textwidth]{screenshots/assistant-1280px.png}} \\[2pt]
 {\scriptsize (c) Hàng đợi duyệt bài chặn người tạo tự duyệt bài của mình} &
 {\scriptsize (d) Màn hình Trợ lý Marketing AI và các mẫu câu lệnh gợi ý}
\end{tabular}
\caption{Ảnh chụp màn hình thực tế các chức năng chính của hệ thống MarketFlow}
\label{fig:ui-wireframe}
\end{figure}

\section{Tích hợp Trợ lý AI và thiết kế System Prompt}

\subsection{Cơ chế truyền quyền tạm thời (\texttt{Run-Token}) và đổi mô hình tự động}

Khi người dùng gửi câu hỏi cho Trợ lý AI, Backend tạo một mã JWT ngắn hạn (\texttt{typ="assistant\_run"}, có hiệu lực 10 phút) chứa \texttt{user\_id}, \texttt{role} cùng \texttt{thread\_id}, sau đó gửi sang dịch vụ \texttt{ai-runner}. Khi tiến trình AI gọi vào cổng \texttt{/mcp}, máy chủ giải mã token này để kiểm tra quyền: nếu người đang chat là \texttt{MARKETER}, AI cũng chỉ được tạo đề xuất trên đúng các chiến dịch mà nhân viên đó phụ trách.

Để hệ thống luôn hoạt động khi dùng các gói API miễn phí, nhóm cấu hình 13 mô hình thuộc 4 nhà cung cấp trong \texttt{deploy/opencode/opencode.json} và \texttt{ModelSelector.tsx} (Bảng~\ref{tab:cost-table}). Nếu mô hình đang chọn bị giới hạn lượt gọi (HTTP 429) hoặc không phản hồi, \texttt{ai-runner} tự động đánh dấu tạm nghỉ mô hình đó 180 giây và chuyển sang mô hình dự phòng kế tiếp.

\begin{table}[!htbp]
\centering
\scriptsize
\renewcommand{\arraystretch}{1.04}
\caption{Danh sách 13 mô hình LLM được cấu hình trong hệ thống}
\label{tab:cost-table}
\begin{tabularx}{\textwidth}{|L{2.6cm}|L{5.5cm}|C{1.8cm}|Y|}
\hline
\thead{Nhà cung cấp} & \thead{Tên mô hình trong cấu hình (13 mô hình)} & \thead{Chi phí} & \thead{Vai trò sử dụng}\\
\hline
\textbf{OpenCode Free} (3 mô hình) & \texttt{opencode/deepseek-v4.1-flash-free}, \texttt{opencode/glm-5.3-flash-free}, \texttt{opencode/mimo-v2.6-flash-free} & Miễn phí & Nhóm mặc định, hỗ trợ gọi chuỗi công cụ MCP ổn định.\\
\hline
\textbf{Groq Cloud} (3 mô hình) & \texttt{groq/openai/gpt-oss-120b}, \texttt{groq/qwen/qwen3.8-27b}, \texttt{groq/openai/gpt-oss-20b} & Miễn phí (Free Tier) & Nhóm tốc độ phản hồi nhanh, dùng khi cần kết quả tức thì.\\
\hline
\textbf{OpenRouter} (5 mô hình) & \texttt{nvidia/nemotron-3.5-super-vl:free}, \texttt{nvidia/nemotron-3-ultra-550b-a55b:free}, \texttt{nvidia/nemotron-3.5-lightning:free}, \texttt{thinkingmachines/inkling:free}, \texttt{inkling-small:free} & Miễn phí (Free Tier) & Nhóm dự phòng khi các cổng chính hết hạn mức theo phút.\\
\hline
\textbf{Google AI Studio} (2 mô hình) & \texttt{google/gemini-3.5-flash-lite}, \texttt{google/gemini-3.1-flash-lite} & Miễn phí (Free Tier) & Nhóm dự phòng xử lý tiếng Việt và ngữ cảnh dài.\\
\hline
\end{tabularx}
\end{table}

\subsection{Thiết kế System Prompt, định dạng Input/Output và cải tiến qua 4 phiên bản}

System Prompt của Trợ lý AI được tách riêng khỏi mã nguồn Python, lưu tại \texttt{deploy/opencode/\allowbreak agent/marketflow.md} và \texttt{deploy/opencode/prompts/}. Bảng~\ref{tab:prompt-io-spec} trích dẫn nội dung thực tế của System Prompt (Bản 4), câu lệnh người dùng mẫu và cấu trúc JSON đầu vào/đầu ra khi gọi công cụ MCP.

\begin{table}[!htbp]
\centering
\scriptsize
\renewcommand{\arraystretch}{1.03}
\caption{Trích dẫn System Prompt thực tế, User Prompt mẫu và định dạng Input/Output}
\label{tab:prompt-io-spec}
\begin{tabularx}{\textwidth}{|L{2.9cm}|Y|}
\hline
\thead{Thành phần} & \thead{Nội dung cấu hình thực tế trong hệ thống}\\
\hline
\textbf{System Prompt}\par (\texttt{deploy/opencode/\allowbreak agent/marketflow.md}) &
\texttt{permission: \{"*": "deny", "external\_directory": "deny", "marketflow\_*": "allow"\}}\par
1. \textbf{Trung thực về dữ liệu}: Chỉ dùng con số có thực từ công cụ \texttt{marketflow\_*} hoặc \texttt{[NGỮ CẢNH]}. Nếu chưa có số liệu, ghi rõ ``chưa có dữ liệu thống kê''. Công thức: $\text{CTR}=\text{Clicks}/\text{Views}$, $\text{CVR}=\text{Conversions}/\text{Clicks}$, $\text{CPC}=\text{Cost}/\text{Clicks}$, $\text{CPA}=\text{Cost}/\text{Conversions}$.\par
2. \textbf{Chỉ tạo đề xuất (\texttt{PROPOSED})}: Mọi thao tác ghi phải gọi \texttt{propose\_*}. Không tự tuyên bố đã duyệt hay đã đăng bài. Sau khi gọi \texttt{propose\_*}, không chép lặp lại toàn bộ bài nháp vào khung chat.\par
3. \textbf{Giới hạn độ dài kênh \& từ cấm}: Facebook $<$ 1.500 ký tự, kịch bản TikTok $<$ 1.200 ký tự (chia 3 phần \texttt{[Hook 0-3s]}, \texttt{[Nội dung 3-25s]}, \texttt{[CTA 25-30s]}), tiêu đề Email $<$ 60 ký tự. Không dùng từ cấm (\textit{``cam kết 100\%'', ``thần dược'', ``rẻ nhất''}).\\
\hline
\textbf{User Prompt mẫu}\par \& Ngữ cảnh đính kèm &
\texttt{[NGỮ CẢNH] campaign\_id=1, name="Ra Mắt Dòng Cold Brew Đóng Chai Tiện Lợi", goal="SALES", role="MARKETER"}\par
\textbf{Câu hỏi}: \textit{``Hãy đề xuất 1 bản nháp bài đăng Facebook cho chiến dịch Cold Brew, bám sát mục tiêu tăng doanh số.''}\\
\hline
\textbf{Đầu vào công cụ}\par (\texttt{MCP Tool Input JSON}) &
\texttt{propose\_content\_drafts(\{"campaign\_id": 1, "rationale": "Tập trung tiện lợi pha sẵn cho dân văn phòng", "items": [\{"channel\_code": "FACEBOOK", "content\_type": "SOCIAL\_POST", "title": "Cold Brew Ủ Lạnh 18 Tiếng", "body": "...", "cta": "Đặt mua combo 3 chai ngay", "hashtags": "\#AnNamCoffee \#ColdBrew"\}]\})}\\
\hline
\textbf{Đầu ra công cụ}\par (\texttt{MCP Tool Output JSON}) &
\texttt{\{"action\_id": 14, "action\_type": "CREATE\_CONTENT\_BATCH", "status": "PROPOSED", "banned\_warnings": [], "message": "Đã tạo thẻ đề xuất \#14 chờ người dùng xác nhận."\}}\\
\hline
\end{tabularx}
\end{table}

Nhóm đã thử nghiệm và chỉnh sửa System Prompt qua 4 phiên bản:
\begin{itemize}[noitemsep,topsep=2pt]
 \item \textbf{Bản 1 (Cơ sở)}: Chỉ mô tả ngắn gọn vai trò trợ lý marketing. Ở bản này, AI đôi khi tự bịa thêm con số giả định khi chiến dịch chưa có dữ liệu hoặc viết kịch bản TikTok quá dài (hơn 3.600 ký tự).
 \item \textbf{Bản 2 (Thêm quy tắc an toàn)}: Bổ sung yêu cầu không bịa số liệu và chỉ tạo đề xuất chờ duyệt. Tuy nhiên, do viết câu dặn dò quá dài, AI lại sinh kèm một đoạn cảnh báo dài dòng ở cuối mỗi bài viết khiến độ dài bài đăng Facebook và TikTok vượt quá giới hạn ký tự cho phép.
 \item \textbf{Bản 3 (Thêm công thức KPI và định dạng kênh)}: Quy định rõ cách đọc chỉ số từ hàm \texttt{get\_stats} và cấu trúc bài viết từng kênh (mở bài, nội dung, CTA, hashtag).
 \item \textbf{Bản 4 (Bản chính thức đang chạy)}: Bỏ các câu rào đón thừa, quy định cứng giới hạn độ dài từng kênh, bắt buộc gọi \texttt{check\_banned\_keywords} trước khi đề xuất bài nháp và từ chối các câu lệnh yêu cầu vượt quyền.
\end{itemize}

\section{Nhật ký sử dụng AI trong quá trình phát triển phần mềm (SDLC)}

Trong quá trình làm đồ án, nhóm sử dụng công cụ lập trình AI để hỗ trợ ở cả 4 giai đoạn (từ thiết kế bảng, viết khung CRUD đến viết test). Thay vì dùng nguyên mã do AI sinh ra, nhóm đã kiểm tra lại và tự sửa 4 lỗi thực tế như ghi nhận trong Bảng~\ref{tab:sdlc-ai-log}. Bảng~\ref{tab:sdlc-code-diff} minh họa cụ thể câu lệnh nhóm đã hỏi AI, đoạn mã bị lỗi do AI sinh ra và đoạn mã nhóm đã tự sửa lại trong dự án.

\begin{table}[!htbp]
\centering
\scriptsize
\renewcommand{\arraystretch}{1.02}
\caption{Nhật ký sử dụng AI trong quá trình làm dự án và các lỗi nhóm đã tự khắc phục}
\label{tab:sdlc-ai-log}
\begin{tabularx}{\textwidth}{|C{1.8cm}|L{2.8cm}|Y|Y|}
\hline
\thead{Giai đoạn} & \thead{Việc nhờ AI hỗ trợ} & \thead{Vấn đề trong kết quả do AI tạo ra} & \thead{Cách nhóm kiểm tra và sửa lại}\\
\hline
\textbf{KT1}\par Phân tích \& Thiết kế & Gợi ý các bảng CSDL và danh sách công cụ MCP. & AI đề xuất cho các hàm MCP ghi thẳng vào bảng \texttt{campaigns} và \texttt{contents}, dễ làm hỏng dữ liệu nếu AI gọi nhầm hàm. & Thêm bảng trung gian \texttt{ai\_actions}: mọi hàm \texttt{propose\_*} chỉ tạo bản ghi \texttt{PROPOSED} chờ người dùng bấm xác nhận.\\
\hline
\textbf{KT2}\par Xây dựng Backend \& UI & Viết khung dịch vụ \texttt{ai-runner}, luồng SSE và kết nối \texttt{/mcp}. & (1) AI dùng chung một khóa tĩnh cho \texttt{/mcp} nên mất phân quyền theo người dùng. (2) Không giới hạn số tiến trình \texttt{opencode} chạy cùng lúc. (3) Trên frontend dùng \texttt{EventSource} mặc định nên không gửi kèm được header \texttt{Authorization}. & (1) Tự viết hàm \texttt{create\_run\_token} (JWT 10 phút chứa \texttt{user\_id, role}). (2) Đo thử trên VPS thấy mỗi tiến trình tốn $\sim$490 MB RAM nên đặt \texttt{Semaphore(1)} chạy tuần tự. (3) Đổi sang dùng \texttt{fetch()} + \texttt{ReadableStream}.\\
\hline
\textbf{KT3}\par Tối ưu Prompt \& Test & Viết System Prompt và kịch bản chấm điểm tự động. & Ở Bản 2, AI thêm quá nhiều câu lưu ý an toàn ở cuối bài viết làm vượt giới hạn ký tự kênh Facebook và TikTok (tỷ lệ đạt giảm từ 6/8 xuống 5/8). & Sửa lại Bản 4: cấm chèn lời giải thích thừa bên ngoài bài nháp và đặt giới hạn ký tự cụ thể cho từng kênh, nâng kết quả lên 7/8 kịch bản.\\
\hline
\textbf{Cuối kỳ}\par Đưa lên VPS & Viết cấu hình Nginx và Systemd trên Ubuntu VPS. & Cấu hình Nginx do AI gợi ý mở toàn bộ đường dẫn của Backend ra ngoài (lộ cổng \texttt{/mcp}) và bật bộ đệm làm nghẽn luồng chữ SSE. & Thêm cấu hình chặn \texttt{location \^{}\~{} /mcp \{ return 404; \}} và tắt bộ đệm \texttt{proxy\_buffering off} cho đường dẫn stream SSE.\\
\hline
\end{tabularx}
\end{table}

\begin{table}[!htbp]
\centering
\scriptsize
\renewcommand{\arraystretch}{1.03}
\caption{Minh chứng cụ thể câu lệnh hỏi AI, mã lỗi AI sinh ra và mã nhóm đã tự sửa lại}
\label{tab:sdlc-code-diff}
\begin{tabularx}{\textwidth}{|L{3.2cm}|Y|Y|}
\hline
\thead{Câu lệnh nhóm hỏi AI} & \thead{Kết quả AI sinh ra (Có lỗ hổng)} & \thead{Mã nhóm đã kiểm tra và sửa lại trong code}\\
\hline
\textit{``Viết hàm MCP tạo bản nháp bài viết cho chiến dịch và xác thực kết nối giữa ai-runner với cổng /mcp''} &
\texttt{MCP\_API\_KEY = "static-secret"}\par
\texttt{def create\_content(campaign\_id, body):}\par
\texttt{~~db.add(Content(campaign\_id=campaign\_id, body=body, status="APPROVED"))}\par
\texttt{~~db.commit()}\par
$\rightarrow$ \textbf{Lỗi}: Dùng khóa tĩnh nên \texttt{marketer2} có thể nhờ AI sửa chiến dịch của \texttt{marketer1}, và AI ghi thẳng bài \texttt{APPROVED} vào bảng \texttt{contents}. &
\texttt{claims = verify\_run\_token(token) \# JWT 10p}\par
\texttt{require\_campaign\_owner(db, campaign\_id, claims.user\_id, claims.role)}\par
\texttt{action = AIAction(thread\_id=claims.thread\_id, action\_type="CREATE\_CONTENT\_BATCH", status="PROPOSED", payload=items)}\par
\texttt{db.add(action); db.commit()}\\
\hline
\textit{``Viết hàm gọi tiến trình opencode run trong ai-runner khi có yêu cầu chat từ người dùng''} &
\texttt{async def run\_endpoint(req: RunRequest):}\par
\texttt{~~proc = await asyncio.create\_subprocess\_exec("opencode", "run", req.prompt)}\par
$\rightarrow$ \textbf{Lỗi}: Khi 3 người bấm chat cùng lúc, 3 tiến trình ngốn $3 \times 490\text{ MB} \approx 1,5\text{ GB RAM}$ làm VPS 2,5 GB bị tràn bộ nhớ (OOM). &
\texttt{\_SEM = asyncio.Semaphore(1) \# Hàng đợi tuần tự}\par
\texttt{async def run\_endpoint(req: RunRequest):}\par
\texttt{~~async with \_SEM:}\par
\texttt{~~~~for model in fallback\_chain: \# Tự đổi model khi lỗi 429}\par
\texttt{~~~~~~proc = await run\_with\_timeout(model, req)}\\
\hline
\end{tabularx}
\end{table}
